\subsection{有限域的结构}

命题 1. --- 设 K 是一个含有 q 个元素的有限域。

\begin{enumerate}
\def\labelenumi{\alph{enumi})}
\item
  \(\mathbf{K}\) 的特征是一个素数 \(p\) ，并且存在一个整数 \(f \geqslant 1\) 使得 \(q = p^f\) 。
\item
  \(\mathbf{K}\) 的加法群是 \(f\) 个 \(p\) 阶循环群的直和。
\item
  \(\mathbf{K}\) 的乘法群是阶为 \(q - 1\) 的循环群。
\end{enumerate}

由于 Z 是无限的而 K 是有限的，唯一的环同态 \(\varphi:Z\to K\) 不是单射，且其核是 Z 的一个非零素理想。因此 K 的特征是一个素数，且 K 是 \(F_{p}\) （ p 个元素的域）上的代数（V，第2页）。设 f 为 K 在 \(F_{p}\) 上的次数。若 f 为无穷，则 K 将对每个整数 \(n\geqslant0\) 包含 \(F_{p}\) 上的一个 n 维子空间，从而 \(q\geqslant p^{n}\) ，这显然是荒谬的。因此 f 是有限的。于是 K 的加法群同构于 \((\mathbf{F}_{p})^{f}\) ，由此得出断言 a) 和 b)。

断言 c) 由引理 1（V，第78页）得出。

命题 2. --- 设 K 为含 q 个元素的有限域。域 K 是多项式 \(X^{q}-X\) （ \(F_{p}[X]\) 中的）的一个分裂域，并且它就是该多项式全部根组成的集合。

对于每一个 \(x \neq 0\) （在 \(\mathbf{K}\) 中），我们有 \(x^{q-1} = 1\) ，因为 \(\mathbf{K}^*\) 是一个有限群，其阶为 \(q - 1\) （I，第52页）。由此可知 \(x^q = x\) 对每个 \(x\) （在 \(\mathbf{K}\) 中）成立。多项式 \(\mathbf{X}^q - \mathbf{X}\) （ \(\mathbf{F}_p[\mathbf{X}]\) 中的）的次数为 \(q\) ，且它有 \(q\) 个根（在 \(\mathbf{K}\) 中），从而

\[
\mathbf {X} ^ {q} - \mathbf {X} = \prod_ {\xi \in K} (\mathbf {X} - \xi).\tag{1}
\]

命题 2 由此立即得出。

推论. --- 两个基数相同的有限域是同构的。

设 \(\mathbf{K}'\) 为一个具有 \(q\) 个元素的有限域；其特征是一个素数 \(p'\) ，它整除 \(q = p^f\) ，从而 \(p' = p\) 。所以 \(\mathbf{K}'\) 是多项式 \(\mathbf{X}^q - \mathbf{X}\) （ \(\mathbf{F}_p[\mathbf{X}]\) 中的）的一个分裂域（命题 2），因此 \(\mathbf{K}\) 与 \(\mathbf{K}'\) 同构（V，第22页，推论）。

当 \(\mathbf{K} = \mathbf{F}_p\) 时，公式（1）化简为关系式

\[
\mathbf {X} ^ {p} - \mathbf {X} \equiv \prod_ {i = 0} ^ {p - 1} (\mathbf {X} - i) \bmod p \mathbf {Z} [ \mathbf {X} ]\tag{2}
\]

在多项式环 Z{[}X{]} 中。

公式（2）也可写成

\[
\mathbf {X} ^ {p - 1} - 1 \equiv \prod_ {i = 1} ^ {p - 1} (\mathbf {X} - i) \bmod p \mathbf {Z} [ \mathbf {X} ].\tag{3}
\]

特别地，当 X = 0 时，我们得到（«威尔逊公式»）

\[
(p - 1)! \equiv - 1 \bmod p.\tag{4}
\]

